\chapter{Marco Teórico / Estado del Arte}

Este capítulo revisa el conocimiento científico relevante para el trabajo, organizado en cuatro
áreas: la generación automática de informes radiológicos, la explicabilidad de los modelos
\textit{transformer}, la estimación de incertidumbre y la detección de alucinaciones en modelos
generativos, y los conjuntos de datos y modelos preentrenados del dominio radiológico. El capítulo
concluye con una síntesis que sitúa la aportación del trabajo respecto al estado del arte y delimita
el hueco de investigación que se aborda.

\section{Generación automática de informes radiológicos}

\subsection{De la descripción de imágenes al informe clínico}
La generación de informes a partir de imágenes médicas hereda su formulación del \textit{image
captioning}: producir una secuencia de texto condicionada a una imagen. Las primeras aproximaciones
adoptaron arquitecturas \textbf{codificador--decodificador} en las que una red convolucional
resumía la imagen en un vector de características y una red recurrente (LSTM o GRU) generaba el texto
palabra a palabra. La incorporación de mecanismos de \textbf{atención visual} permitió que cada
palabra generada se condicionara a distintas regiones de la imagen, mejorando tanto la calidad como
la interpretabilidad aparente de la salida.

El informe radiológico presenta, sin embargo, particularidades que lo distinguen de una simple
descripción. Es un texto \textbf{largo, estructurado y altamente estandarizado}, organizado en
secciones (\textit{Findings}, \textit{Impression}) y dominado por un lenguaje clínico repetitivo en
el que la mayoría de estudios son normales. Estas características motivaron arquitecturas específicas:
modelos con \textbf{memoria relacional} y módulos de generación guiada por conocimiento, que buscaban
capturar la estructura del informe y las correlaciones entre hallazgos, y que constituyeron el
estándar antes de la irrupción de los grandes modelos.

\subsection{La era de los modelos multimodales de lenguaje}
La generalización de los \textit{transformers} y el preentrenamiento a gran escala dieron lugar a los
\textbf{modelos multimodales de lenguaje} (MLLM), que acoplan un codificador visual a un modelo de
lenguaje de gran tamaño mediante una capa de proyección, siguiendo el paradigma de LLaVA
\parencite{liu2023llava}. Su adaptación al dominio radiológico ha producido generadores de informes
de calidad notablemente superior. Entre ellos destacan \textbf{LLaVA-Rad} \parencite{llavarad2024},
que combina un codificador biomédico con un decodificador Vicuna afinado sobre datos de tórax; y
\textbf{MAIRA-2} \parencite{maira2}, orientado a la \textit{generación anclada} (\textit{grounded
report generation}): además del texto, produce \textbf{cajas delimitadoras} que localizan cada
hallazgo, entrenado sobre una combinación de corpus (MIMIC-CXR, PadChest y USMix). Otros sistemas
recientes exploran el razonamiento por agentes y la integración de herramientas de detección
especializadas. La Tabla~\ref{tab:modelos} resume las características de los modelos empleados o
relevantes para este trabajo.

\begin{table}[t]\centering
\caption{Modelos multimodales de CXR relevantes. RRG = \textit{radiology report generation}.}
\label{tab:modelos}
\begin{tabular}{lllc}
\toprule
Modelo & Base de lenguaje & Entrenamiento & Anclaje explícito \\ \midrule
LLaVA-Rad & Vicuna-7B & MIMIC-CXR & No \\
MAIRA-2   & Vicuna-7B & MIMIC-CXR, PadChest, USMix & Sí (cajas) \\
CXR-LLAVA & Llama-2-7B & Multi-corpus & No \\ \bottomrule
\end{tabular}
\end{table}

\subsection{El problema de la evaluación}
La evaluación de estos sistemas se apoya tradicionalmente en \textbf{métricas de solapamiento léxico}
heredadas del resumen automático, como ROUGE \parencite{lin2004rouge}, BLEU o METEOR, que comparan el
texto generado con el informe de referencia contando $n$-gramas comunes. Estas métricas presentan una
debilidad grave en radiología: al estar el dominio dominado por informes normales y por un lenguaje
de plantilla, \textbf{un modelo que memoriza la plantilla mayoritaria obtiene puntuaciones altas sin
leer realmente la imagen}. Este trabajo constata empíricamente el fenómeno: una variante mejor
anclada a la imagen puede obtener un ROUGE inferior al de un modelo ciego que reproduce la plantilla.

Por ello, la comunidad ha desarrollado \textbf{métricas de eficacia clínica} que evalúan la
corrección de los hallazgos y no la coincidencia superficial de palabras: clasificadores de
etiquetas radiológicas que extraen los hallazgos de ambos informes y los comparan (obteniendo
precisión, sensibilidad y $F_1$ por patología), y métricas basadas en grafos de conocimiento que
contrastan las entidades clínicas y sus relaciones. Entre las primeras destacan los
\textbf{clasificadores de etiquetas radiológicas} como \textbf{CheXbert} \parencite{smit2020chexbert},
que combina reglas expertas y un modelo BERT afinado para extraer de cada informe un vector de
hallazgos —presente, ausente o incierto para cada patología— y calcular métricas de precisión,
sensibilidad y $F_1$ sobre esa comparación; entre las segundas, las basadas en \textbf{grafos de
conocimiento} como \textbf{RadGraph} \parencite{jain2021radgraph}, que representa las entidades
clínicas de un informe (hallazgos, anatomía, atributos) y sus relaciones en un grafo estructurado, y
cuya comparación entre grafo generado y grafo de referencia (\textit{RadGraph~$F_1$}) penaliza tanto
los hallazgos inventados como los omitidos con mayor granularidad que un simple recuento de
etiquetas. Estas métricas se aproximan mucho más al criterio del radiólogo que el solapamiento léxico,
aunque a costa de depender de un extractor de hallazgos que introduce su propio error —una limitación
que este trabajo evita adoptando, en su lugar, fuentes de verdad ya estructuradas (MeSH, hallazgos de
radiólogo) y un juez de lenguaje sobre la narrativa completa (Capítulo~4)—. La distinción
entre \emph{corrección clínica} y
\emph{similitud textual} es central para este trabajo, pues motiva que la fiabilidad se defina en
términos de acierto sobre hallazgos significativos y no de solapamiento léxico. En consecuencia, la
frontera de la investigación se ha desplazado de \emph{generar mejor} a \emph{saber cuándo fiarse} de
lo generado: la fiabilidad y la explicabilidad.

\section{Anatomía y funcionamiento de un modelo multimodal de lenguaje}
Comprender cómo funciona internamente un MLLM es imprescindible para justificar por qué las señales
empleadas en este trabajo —el mapa de calor de atención y las probabilidades de los tokens— son
precisamente las que son. Esta sección describe sus componentes.

\subsection{El \textit{transformer} y el mecanismo de atención}
La arquitectura \textit{transformer} \parencite{vaswani2017} sustituyó la recurrencia por el
mecanismo de \textbf{auto-atención}: cada elemento de una secuencia calcula una combinación ponderada
del resto, donde los pesos —la \emph{atención}— se obtienen de la similitud entre una \textit{query} y
unas \textit{keys}. Formalmente, $\mathrm{Attention}(Q,K,V)=\mathrm{softmax}(QK^{\top}/\sqrt{d})V$.
Este mecanismo, además de su eficacia, expone una magnitud interpretable —cuánto atiende cada token a
cada otro— que está en el origen del debate sobre si la atención constituye una explicación, y que
este trabajo aprovecha y, al mismo tiempo, somete a validación.

\subsection{El codificador visual}
Para procesar la imagen, los MLLM emplean un \textbf{Vision Transformer} (ViT)
\parencite{dosovitskiy2021}, que divide la imagen en una rejilla de \textbf{parches} (por ejemplo,
$16\times16$ o $14\times14$ píxeles), los proyecta linealmente en \textit{tokens} y les aplica
auto-atención como si fueran palabras. El resultado es una secuencia de \textbf{características por
parche} más un \textit{token} global (CLS). Esta rejilla de parches es la que confiere \emph{estructura
espacial} a la atención: al promediar la atención que los tokens de texto dirigen a los tokens
visuales y reordenarla en la rejilla, se obtiene el mapa de calor sobre la imagen. En los codificadores
empleados en este trabajo, dicha rejilla es de $37\times37$ parches. Codificadores auto-supervisados de
dominio, como RAD-DINO \parencite{perezgarcia2024}, aprenden estas representaciones sin supervisión
textual y mejoran la calidad del anclaje.

\subsection{El conector y el preentrenamiento contrastivo}
El puente entre la imagen y el modelo de lenguaje —el \textbf{conector}— adopta distintas formas: una
simple \textbf{proyección lineal} de las características visuales al espacio del lenguaje (como en
LLaVA \parencite{liu2023llava}), un módulo de consulta aprendido (\textit{Q-Former}) que resume la
imagen en unos pocos tokens (BLIP-2 \parencite{li2023blip2}), o \textbf{atención cruzada con
compuerta} intercalada en el decodificador (Flamingo \parencite{alayrac2022}). El generador propio de
este trabajo adopta esta última vía. Con frecuencia, el codificador visual se ha alineado previamente
con el texto mediante un \textbf{preentrenamiento contrastivo} de tipo CLIP \parencite{radford2021},
que aproxima en un espacio común los embeddings de imágenes y textos emparejados; la Etapa~1 del
modelo propio reproduce esta idea.

\subsection{El decodificador y la generación}
Una vez inyectadas las características visuales, un \textbf{modelo de lenguaje} decodificador genera el
informe de forma \textbf{autorregresiva}: en cada paso produce una distribución de probabilidad sobre
el vocabulario y emite un token. Este proceso es la segunda fuente de señal que explota el trabajo: la
\textbf{probabilidad} que el modelo asigna a cada token generado —y su entropía— son un subproducto
directo del paso hacia adelante, disponible \emph{sin coste adicional} en cualquier generador. El
entrenamiento suele constar de una fase de alineamiento y una de \textit{instruction tuning} visual
\parencite{liu2023llava}, que enseña al modelo a seguir la instrucción de describir la imagen.

\paragraph{Consecuencia metodológica.} La arquitectura anterior explica por qué las dos señales
centrales del trabajo son \emph{naturales} y \emph{agnósticas al modelo}: la atención de los tokens de
texto hacia los tokens visuales proporciona el \textbf{mapa de calor} (la explicación de \emph{dónde}
mira el modelo), y la distribución de probabilidad de cada token proporciona la \textbf{incertidumbre}
(la señal de \emph{cuánto} fiarse). Ambas se obtienen del mismo paso de inferencia que ya se ejecuta
para generar el informe, sin modelos auxiliares ni cómputo extra.

\section{Por qué abordar el problema desde la fiabilidad y la incertidumbre}
La decisión de centrar el trabajo en la fiabilidad —y no en mejorar la generación— responde a un
diagnóstico del estado del arte. Por un lado, la calidad del texto generado por los MLLM de CXR
\emph{se está saturando}, mientras que el verdadero cuello de botella para su adopción clínica es la
\textbf{confianza}: la tendencia de estos modelos a alucinar sigue siendo la principal barrera
\parencite{ji2023}. Por otro, existen dos grandes familias de señales de fiabilidad: las
\textbf{externas}, que requieren un verificador adicional (un informe de referencia u otro modelo), y
las \textbf{intrínsecas}, que se leen de la propia salida del modelo. Este trabajo prioriza las
segundas por tres razones: son \textbf{agnósticas al modelo} (transfieren a cualquier generador), están
\textbf{disponibles en inferencia} sin coste adicional, y capturan la información de que el propio
modelo dispone sobre su acierto —los modelos de lenguaje, en buena medida, ``saben lo que saben''
\parencite{kadavath2022}—. La cuestión abierta, que el trabajo aborda, es si esa información se lee
mejor en la \emph{forma} de la explicación visual o en la \emph{incertidumbre} de la decisión, y si,
una vez elegida, resulta \textbf{calibrable} \parencite{guo2017} para su uso clínico.

\section{Explicabilidad en modelos \textit{transformer}}

\subsection{¿Es la atención una explicación?}
La explicabilidad de los \textit{transformers} gira en gran medida en torno a si los pesos de
\textbf{atención} constituyen una explicación válida de sus decisiones. El debate quedó planteado por
\textcite{jain2019}, que mostraron que es posible encontrar distribuciones de atención muy distintas
que producen la misma predicción, cuestionando que la atención sea \emph{fiel} al razonamiento del
modelo. La réplica de \textcite{wiegreffe2019} matizó la conclusión: la atención sí porta información
explicativa bajo ciertas condiciones y su utilidad depende de cómo se defina y evalúe una
explicación. Este debate es directamente pertinente para la radiología, donde el mapa de calor de
atención se ofrece con frecuencia como justificación visual del diagnóstico.

\subsection{Métodos de atribución y agregación}
Para agregar la atención a lo largo de las múltiples capas del modelo se propuso el \textit{attention
rollout} \parencite{abnar2020}, que compone las matrices de atención capa a capa teniendo en cuenta
las conexiones residuales. Métodos posteriores como el de \textcite{chefer2021} combinan atención y
gradientes para obtener mapas de relevancia más fieles, propagando la relevancia a través de la red
de forma consistente con las reglas de conservación. En paralelo, la familia de métodos basados en
\textbf{gradientes} produce mapas de saliencia a partir de los gradientes de la salida respecto a las
activaciones intermedias: el \textit{class activation mapping} y su generalización \textbf{Grad-CAM}
ponderan los mapas de características por la importancia del gradiente para obtener una localización
gruesa de la evidencia, mientras que los métodos de \textbf{propagación de relevancia} redistribuyen la
puntuación de salida hacia atrás, capa a capa, hasta la entrada. Cada familia tiene ventajas e
inconvenientes: los mapas de atención son baratos pero de fidelidad discutida; los basados en
gradientes son más costosos y sensibles al ruido; y los de propagación exigen reglas específicas por
tipo de capa. Frente a ellos, \textit{Integrated
Gradients} \parencite{sundararajan2017} mide la contribución causal de cada entrada integrando los
gradientes a lo largo de un camino desde una referencia neutra hasta la entrada real, ofreciendo
\textbf{garantías axiomáticas} (sensibilidad e invariancia de implementación) de las que la atención
carece.

\subsection{Evaluación objetiva de la fidelidad}
La proliferación de métodos hace imprescindible \textbf{evaluar objetivamente} la fidelidad de una
explicación, en lugar de limitarse a comparar mapas entre sí. Las \textbf{métricas de perturbación}
—eliminación e inserción de las regiones más relevantes— y las nociones de
\textit{comprehensiveness} (cuánto cae la confianza al eliminar lo señalado como importante) y
\textit{sufficiency} (cuánta confianza se conserva reteniendo solo lo importante), formalizadas en el
\textit{benchmark} ERASER \parencite{deyoung2020}, permiten cuantificar si la explicación captura de
verdad lo que impulsa la predicción. Este trabajo adopta esta filosofía: la explicación se valida
frente a \textbf{referencias objetivas} —cajas delimitadoras reales y perturbaciones controladas— en
lugar de asumir su fidelidad por inspección visual. Como advertencia transversal, \textcite{ghassemi2021}
argumentan que las explicaciones locales por saliencia en salud pueden generar una
\textbf{falsa sensación de confianza} y fomentar el sesgo de automatización si se presentan como
garantía de fiabilidad sin validación, lo que refuerza la necesidad de distinguir entre \emph{dónde}
mira el modelo y \emph{cuánto} debe fiarse el clínico de su veredicto.

\section{Incertidumbre y detección de alucinaciones}

\subsection{Naturaleza de las alucinaciones}
En los modelos generativos de lenguaje se denomina \textbf{alucinación} a la producción de
afirmaciones no sustentadas por la evidencia. Los estudios de síntesis del fenómeno
\parencite{ji2023} distinguen entre alucinaciones \emph{intrínsecas} (que contradicen la entrada) y
\emph{extrínsecas} (no verificables con ella), y las atribuyen a factores tanto de los datos —sesgos y
ruido en el corpus de entrenamiento— como de la propia decodificación —la exposición a su propia
salida y la tendencia a completar patrones plausibles—. En
la generación de informes radiológicos, esta taxonomía se concreta en dos modos de error
clínicamente opuestos: la \textbf{alucinación} propiamente dicha —afirmar un hallazgo agudo que no
está presente (falso positivo)— y la \textbf{omisión} o \textit{miss} —no reportar un hallazgo agudo
que sí lo está (falso negativo)—. Ambos comprometen la seguridad del paciente de forma distinta y,
como se mostrará, presentan una \textbf{firma de incertidumbre radicalmente diferente}, lo que
constituye uno de los ejes del presente trabajo.

\subsection{Cuantificación de la incertidumbre}
La detección de alucinaciones se ha vinculado más a la \textbf{incertidumbre del propio modelo} que a
sus explicaciones visuales. Las señales más directas son de \textbf{nivel de token}: la probabilidad
asignada por el modelo a cada token generado y la \textbf{entropía predictiva} de la distribución en
cada paso. La probabilidad mínima a lo largo de la secuencia (el token menos seguro) y la fracción de
tokens de baja probabilidad son resúmenes habituales. Estas señales presentan, no obstante,
\textbf{confusores} conocidos —la longitud del texto, la posición del token o su función sintáctica—
que este trabajo controla explícitamente. Un refinamiento propuesto es la \textbf{entropía semántica}
\parencite{kuhn2023}, que agrupa las generaciones equivalentes en significado antes de medir la
incertidumbre, de modo que la variabilidad puramente léxica no infle la estimación.

\subsection{Auto-consistencia y calibración}
Una vía complementaria y \textbf{agnóstica al modelo} (solo requiere poder muestrear) es la
\textbf{auto-consistencia}: SelfCheckGPT \parencite{manakul2023} detecta alucinaciones sin recursos
externos generando varias muestras y midiendo su acuerdo, bajo la hipótesis de que un contenido
factual tiende a repetirse mientras que uno alucinado varía entre muestras. Esta idea se ha adaptado
específicamente a la generación de informes radiológicos mediante métodos de tipo caja negra que
muestrean varios informes y contrastan la consistencia de cada afirmación, como RadFlag
\parencite{radflag2024}, que combina muestreo y un juez de lenguaje para marcar las afirmaciones poco
respaldadas. El presente trabajo se relaciona con esta línea pero se diferencia de ella en que, además
de la consistencia, estudia la señal de probabilidad de token —más barata, pues no exige múltiples
generaciones—, la caracteriza frente a las omisiones y la valida de forma cruzada entre modelos y
dominios. Finalmente, para que una
señal de confianza sea útil no basta con que discrimine: debe estar \textbf{calibrada}, es decir, que
una confianza estimada del 20\,\% corresponda a un 20\,\% de aciertos observados. La calibración se
evalúa con curvas de fiabilidad y con el error de calibración esperado, y se corrige con métodos como
la regresión de Platt o la regresión isotónica. La literatura advierte, además, de que el ajuste por
refuerzo con retroalimentación humana tiende a producir modelos \textbf{sobre-confiados}, lo que
puede degradar precisamente las señales de probabilidad en modelos muy afinados. Estas ideas motivan
que en este trabajo se evalúen la probabilidad de los tokens y la variabilidad entre muestras como
señales candidatas, que se controle su calibración y que se declaren sus límites. La
Tabla~\ref{tab:senales_lit} resume las familias de señales revisadas y su coste.

\begin{table}[t]\centering
\caption{Familias de señales de incertidumbre revisadas en la literatura. Coste en inferencia: sin
generaciones adicionales (bajo) o con varias muestras extra (alto).}
\label{tab:senales_lit}
\begin{tabular}{p{4.3cm}lp{2.9cm}p{3.3cm}}
\toprule
Señal & Nivel & Coste & Referencia \\ \midrule
Probabilidad de token (mín./media) & Token & Bajo & Este trabajo \\
Entropía predictiva & Token & Bajo & --- \\
Entropía semántica & Secuencia & Alto (agrupa muestras) & \parencite{kuhn2023} \\
Auto-consistencia (SelfCheckGPT) & Secuencia & Alto ($K$ muestras) & \parencite{manakul2023} \\
Auto-consistencia + juez (RadFlag) & Secuencia & Alto ($K$ muestras + juez) & \parencite{radflag2024} \\
\bottomrule
\end{tabular}
\end{table}

\section{Datos y modelos preentrenados en radiología}

\subsection{Conjuntos de datos}
La disponibilidad de datos condiciona fuertemente el dominio. El conjunto \textbf{NIH ChestX-ray8}
\parencite{wang2017} aporta clasificación multi-etiqueta y, de forma parcial, \textbf{localización
mediante cajas delimitadoras} para ocho patologías, lo que lo hace idóneo como referencia objetiva de
anclaje. \textbf{Open-I / Indiana University X-ray} \parencite{demnerfushman2016} es prácticamente el
único conjunto de acceso libre con \textbf{informes narrativos completos} y, además, con diagnóstico
estructurado (términos MeSH), lo que permite un etiquetado objetivo del error. Corpus de mayor tamaño
como \textbf{MIMIC-CXR} ofrecen cientos de miles de informes pero requieren credenciales
restringidas. \textbf{PadChest} y su subconjunto anclado \textbf{PadChest-GR} \parencite{padchestgr2024}
aportan un dominio distinto (población española), anotaciones de hallazgos \textbf{presentes y
ausentes} por radiólogo y cajas, resultando idóneos tanto como segunda fuente de verdad objetiva
—independiente del juez basado en modelo de lenguaje— como para evaluar la generalización fuera de
distribución.

\subsection{Codificadores y destilación de conocimiento}
La restricción de datos narrativos motiva el uso de \textbf{destilación de conocimiento} desde
modelos preentrenados libres. En el eje textual, \textbf{CXR-BERT} \parencite{boecking2022} es un
codificador de lenguaje con semántica clínica, entrenado para alinear texto e imagen radiológicos; en
el eje visual, \textbf{RAD-DINO} \parencite{perezgarcia2024} es un codificador de imagen radiológica
auto-supervisado que aprende representaciones potentes sin depender de supervisión textual, y que
constituye, de hecho, el codificador visual de MAIRA-2. La destilación permite mejorar un modelo
propio de escala reducida sin acceso a datos adicionales, convirtiendo la limitación de datos en una
decisión metodológica deliberada. El panorama de generadores del estado del arte es, además,
crecientemente diverso: junto a los modelos densos de tipo LLaVa han aparecido sistemas de
\textit{grounded reporting} que producen cajas, modelos de propósito médico general y enfoques de
razonamiento por agentes que combinan varios modelos especializados. Esta diversidad refuerza el
interés de disponer de una señal de fiabilidad \textbf{agnóstica al generador}, que pueda aplicarse a
cualquiera de ellos sin depender de su arquitectura interna, como la que este trabajo propone.

\section{Síntesis y hueco de investigación}
De la revisión anterior se desprende un panorama consistente. Los MLLM de CXR del estado del arte
generan texto de calidad y muestran alta especificidad, pero adolecen de dos carencias: una
\textbf{sensibilidad limitada} —tienden a omitir hallazgos, con una especificidad notablemente
superior a su sensibilidad ante la patología \parencite{frontiers2025}— y, sobre todo, una
\textbf{ausencia de explicabilidad validada}, como se ha discutido en las Secciones~3.2 y~3.4. Al
mismo tiempo, la literatura de explicabilidad
advierte de que las explicaciones visuales pueden ser \textbf{engañosas} si se asumen fiables sin
validación \parencite{ghassemi2021}, y la literatura de incertidumbre ofrece señales prometedoras
—probabilidad de token, entropía semántica, auto-consistencia— pero desarrolladas, en su mayoría,
sobre modelos de lenguaje generales y sin una validación cruzada entre modelos y dominios en el
ámbito radiológico. Este trabajo se sitúa precisamente en ese hueco: en lugar de proponer otro
generador, \textbf{audita objetivamente} si la explicación de un MLLM de CXR está anclada, comprueba
si su morfología permite detectar los errores, y determina dónde reside realmente la señal de
fiabilidad, validándola de forma cruzada sobre dos modelos y dos dominios y traduciéndola a una
confianza calibrada por informe.
